%------------------------------------------------------------------------------%
\documentclass[fleqn,utf8,aspectratio=169,12pt]{beamer}

\usepackage{calculo_varias_variaveis-1}

\title{Regiões do Plano ou do Espaço}

%------------------------------------------------------------------------------%
\begin{document}

\addtitle

%------------------------------------------------------------------------------%
\section{Intervalos}

%------------------------------------------------------------------------------%
\begin{frame}{Intervalos}
  Subconjuntos contínuos de $\R$
  \bigskip

  \pause
  Intervalo fechado
  \[
    [0, 1] = \left\{ x\in\R\st 0 \leq x \leq 1 \right\}
  \]
  \pause
  Intervalo aberto
  \[
    (-2, 2) = \left\{ x\in\R\st -2 < x < 2 \right\}
  \]
  \pause
  Aberto a esquerda e fechado a direita
  \[
    (1, 5] = \left\{ x\in\R\st 1 < x \leq 5 \right\}
  \]
\end{frame}

%------------------------------------------------------------------------------%
\section{Regiões}

%------------------------------------------------------------------------------%
\begin{frame}{Ponto Interior}
\begin{columns}
\column{0.5\linewidth}
  Um ponto $(a, b)$ em uma região (conjunto) $R$ no plano $xy$ é um
  \emph{ponto interior} de $R$

  \pause
  \medskip
  se é o centro de um disco de raio positivo
  que está inteiramente em $R$
\column{0.5\linewidth}
  \pause
  \includegraphics{ponto_interior}
\end{columns}
\end{frame}

%------------------------------------------------------------------------------%
\begin{frame}{Ponto de Fronteira}
\begin{columns}
\column{0.5\linewidth}
  Um ponto $(a, b)$ é um \emph{ponto de fronteira} de $R$

  \pause
  \medskip
  se todo disco centrado em $(a, b)$ contém ao mesmo tempo pontos
  que estão em $R$ e fora de $R$

  \pause
  \medskip
  O ponto de fronteira não precisa pertencer a $R$
\column{0.5\linewidth}
  \pause
  \includegraphics{ponto_fronteira}
\end{columns}
\end{frame}

%------------------------------------------------------------------------------%
\begin{frame}{Interior de uma Região}
\begin{columns}
\column{0.5\linewidth}
  O conjunto de todos os pontos interiores da região
  formam o \emph{interior} da região
\column{0.5\linewidth}
  \pause
  \includegraphics{interior}
\end{columns}
\end{frame}

%------------------------------------------------------------------------------%
\begin{frame}{Fronteira de uma Região}
\begin{columns}
\column{0.5\linewidth}
  Os pontos de fronteira da região formam sua \emph{fronteira}
\column{0.5\linewidth}
  \pause
  \includegraphics{fronteira}
\end{columns}
\end{frame}

%------------------------------------------------------------------------------%
\begin{frame}{Regiões Abertas e Fechadas}
\begin{columns}[b]
\column{0.5\linewidth}
  Uma região é \emph{aberta} se consiste
  apenas de pontos interiores
  \pause
  \includegraphics{regiao_aberta}
\column{0.5\linewidth}
  \pause
  Uma região é \emph{fechada} se contém
  todos os seus pontos de fronteira
  \pause
  \includegraphics{regiao_fechada}
\end{columns}
\end{frame}

%------------------------------------------------------------------------------%
\begin{frame}{Regiões Limitadas}
\begin{columns}[b]
\column{0.5\linewidth}
  Uma região do plano é limitada se está dentro de um disco de raio fixo
  \pause
  \includegraphics{regiao_limitada}
\column{0.5\linewidth}
  \pause
  Caso contrário ela é ilimitada \\
  \pause
  \includegraphics{regiao_ilimitada}
\end{columns}
\end{frame}

%------------------------------------------------------------------------------%
\section{Exemplos}

%------------------------------------------------------------------------------%
\begin{question}
  Descreva o domínio da função
  \(
    f(x, y) = \sqrt{y-x^2}
  \)
  \begin{enumerate}
    \item Qual a condição que os pontos do domínio deve satisfazer?
    \item Que região é essa? Represente-a graficamente.
    \item Qual o interior da região?
    \item Qual a fronteira da região?
    \item A região é aberta ou fechada?
    \item A região é limitada?
  \end{enumerate}
\end{question}

%------------------------------------------------------------------------------%
\begin{resolution}{Solução}
  \onslide<+->{Condição que deve ser satisfeita para um ponto estar no domínio}
  \begin{align*}
    \onslide<+->{y-x^2 & \geq 0 \\}
    \onslide<+->{y & \geq x^2}
  \end{align*}

  \onslide<+->{
  Domínio
  \[
    D_f = \left\{ (x, y)\in\R^2 \st y \geq x^2\right\}
  \]}
\end{resolution}

%------------------------------------------------------------------------------%
\begin{resolution}{Representação Gráfica}
  \centering
  \includegraphics{ex_dominio_sqrt_y_x2}
\end{resolution}

%------------------------------------------------------------------------------%
\begin{resolution}{Solução}
  \begin{enumerate}
    \item \onslide<+->{Interior da região}
    \onslide<+->{
    \[
      \text{Interior } = \left\{ (x, y)\in\R^2 \st y > x^2\right\}
    \]}
    \vspace{-\baselineskip}
    \item \onslide<+->{Fronteira da região}
    \onslide<+->{
    \[
      \text{Fronteira } = \left\{ (x, y)\in\R^2 \st y = x^2\right\}
    \]}
    \vspace{-\baselineskip}
    \item \onslide<+->A região é aberta ou fechada? \\
    \onslide<+->{Fechada,}
    \onslide<+->{pois \(y = x^2\) pertence ao domínio}
    \item \onslide<+->{A região é limitada?} \\
    \onslide<+->{Ilimitada,}
    \onslide<+->{pois dado qualquer disco sempre teremos um ponto fora dele}
  \end{enumerate}
\end{resolution}

%------------------------------------------------------------------------------%
%------------------------------------------------------------------------------%
\begin{question}
  Descreva o domínio da função
  \(
    f(x, y) = \ln\left(x^2 + y^2 - 4\right)
  \)
  \begin{enumerate}
    \item Qual a condição que os pontos do domínio deve satisfazer?
    \item Que região é essa? Represente-a graficamente.
    \item Qual o interior da região?
    \item Qual a fronteira da região?
    \item A região é aberta ou fechada?
    \item A região é limitada?
  \end{enumerate}
\end{question}

%------------------------------------------------------------------------------%
\begin{resolution}{Solução}
  \onslide<+->{Condição}
  \begin{align*}
    \onslide<+->{x^2 + y^2 - 4 & > 0 \\}
    \onslide<+->{x^2 + y^2 & > 4}
  \end{align*}

  \onslide<+->{Pontos exteriores ao circulo centrado na origem e raio 2}

  \onslide<+->{
  Domínio
  \[
    D_f = \left\{ (x, y)\in\R^2 \st x^2 + y^2 > 4\right\}
  \]}
\end{resolution}

%------------------------------------------------------------------------------%
\begin{resolution}{Representação Gráfica}
  \centering
  \includegraphics{ex_dominio_ln_x2_y2_4}
\end{resolution}

%------------------------------------------------------------------------------%
\begin{resolution}{Solução}
  \begin{enumerate}
    \item \onslide<+->{Interior da região}
    \onslide<+->{
    \[
      \text{Interior } = \left\{ (x, y)\in\R^2 \st x^2 + y^2 > 4\right\}
    \]}
    \vspace{-\baselineskip}
    \item \onslide<+->{Fronteira da região}
    \onslide<+->{
    \[
      \text{Fronteira } = \left\{ (x, y)\in\R^2 \st x^2 + y^2 = 4\right\}
    \]}
    \vspace{-\baselineskip}
    \item \onslide<+->A região é aberta ou fechada? \\
    \onslide<+->{Aberta,}
    \onslide<+->{todos os pontos são pontos interiores}
    \item \onslide<+->{A região é limitada?} \\
    \onslide<+->{Ilimitada,}
    \onslide<+->{pois dado qualquer disco sempre teremos um ponto fora dele}
  \end{enumerate}
\end{resolution}

%------------------------------------------------------------------------------%
%------------------------------------------------------------------------------%
\begin{question}
  Descreva o domínio da função
  \(
    f(x, y) = \frac{\sqrt{x^2-4\,}}{\;\ln(y-1)\;}
  \)
  \begin{enumerate}
    \item Qual a condição que os pontos do domínio deve satisfazer?
    \item Que região é essa? Represente-a graficamente.
    \item Qual o interior da região?
    \item Qual a fronteira da região?
    \item A região é aberta ou fechada?
    \item A região é limitada?
  \end{enumerate}
\end{question}

%------------------------------------------------------------------------------%
\begin{resolution}{Solução}
  Condições para avaliar a função
  \(f(x, y) = \frac{\sqrt{x^2-4\,}}{\;\ln(y-1)\;}\)

  \pause
  \begin{enumerate}[<+->]
    \item \(x^2-4 \geq 0\)
    \item \(\ln(y-1) \neq 0\)
    \item \(y-1 > 0\)
  \end{enumerate}
\end{resolution}

%------------------------------------------------------------------------------%
\begin{resolution}{Solução}
\begin{columns}[t]
\column{0.3\linewidth}
  \onslide<+->{Condição 1}
  \begin{align*}
    \onslide<+->{x^2-4     & \geq 0                       \\}
    \onslide<+->{x^2       & \geq 4                       \\}
    \onslide<+->{\abs{x}   & \geq 2                       \\}
    \onslide<+->{x \leq -2 & \quad\text{ou}\quad x \geq 2}
  \end{align*}
\column{0.3\linewidth}
  \onslide<+->{Condição 2}
  \begin{align*}
    \onslide<+->{\ln(y-1) & \neq 0 \\}
    \onslide<+->{y-1      & \neq 1 \\}
    \onslide<+->{y        & \neq 2}
  \end{align*}
\column{0.3\linewidth}
  \onslide<+->{Condição 3}
  \begin{align*}
    \onslide<+->{y-1 & > 0 \\}
    \onslide<+->{y   & > 1}
  \end{align*}
\end{columns}
\end{resolution}

%------------------------------------------------------------------------------%
\begin{resolution}{Solução}
  Domínio
  \[
    D_f = \left\{ (x, y)\in\R^2 \st
      \abs{x} \geq 2 \;\text{ e }\;
      y > 1          \;\text{ e }\;
      y \neq 2
    \right\}
  \]
\end{resolution}

%------------------------------------------------------------------------------%
\begin{resolution}{Representação Gráfica}
  \centering
  \includegraphics{ex_dominio_sqrt_x2_4_ln_y_1}
\end{resolution}

%------------------------------------------------------------------------------%
\begin{resolution}{Solução}
  \begin{enumerate}
    \item \onslide<+->{Interior da região}
    \onslide<+->{
    \[
      \text{Interior } = \left\{ (x, y)\in\R^2 \st
          \abs{x} > 2 \;\text{ e }\;
          y > 1       \;\text{ e }\;
          y \neq 2
        \right\}
    \]}
    \vspace{-\baselineskip}
    \item \onslide<+->{Fronteira da região}
    \onslide<+->{
    \begin{align*}
      \text{Fronteira } = \Big\{ (x, y)\in\R^2 \st
          & y = 1 \text{ e } \abs{x} \geq 2, \text{ ou }\\[-2mm]
          & y = 2 \text{ e } \abs{x} \geq 2, \text{ ou }
          x = \pm 2 \text{ e } y   \geq 1
        \Big\}
    \end{align*}}
    \vspace{-\baselineskip}
    \item \onslide<+->A região é aberta ou fechada? \\
    \onslide<+->{Nem aberta nem fechada}
    \item \onslide<+->{A região é limitada?} \\
    \onslide<+->{Ilimitada,}
    \onslide<+->{pois dado qualquer disco sempre teremos um ponto fora dele}
  \end{enumerate}
\end{resolution}
%------------------------------------------------------------------------------%

%------------------------------------------------------------------------------%
\section{Lista Mínima}

%------------------------------------------------------------------------------%
\begin{frame}{Lista Mínima}
  Cálculo Vol. 2 do Thomas 12\textsuperscript{a} ed. -- Seção 14.1
  \begin{enumerate}
    \item Estudar o texto da seção
    \item Resolver os exercícios: 1-2, 5-9
  \end{enumerate}

  \vfill
  Atenção: A prova é baseada no livro, não nas apresentações
\end{frame}

%------------------------------------------------------------------------------%
\end{document}
%------------------------------------------------------------------------------%
